% Part: first-order-logic
% Chapter: completeness
% Section: henkin-expansion

\documentclass[../../../include/open-logic-section]{subfiles}

\begin{document}

\olfileid{fol}{com}{hen}

\olsection{توسيع هنكن}

\begin{explain}
من مطالب برهان الاكتمال أن النموذج المبني من مجموعة متسقة
كاملة~$\Gamma$ يرضي ال!!{formula}s المكممة كلها في~$\Gamma$.
ولتحصيل هذا الضمان نأخذ بطريقة ليون هنكن: نزيد اللغة
!!{constant}s لا نهاية لها، ونلحق، لكل !!{formula} لها
!!{variable} حر واحد $!A(x)$، صيغة على النحو
\iftag{prvEx}
      {$\lexists[x][!A(x)] \lif !A(c)$}
      {$\lnot\lforall[x][!A(x)] \lif \lnot !A(c)$},
على أن $c$ من ال!!{constant}s المستحدثة. فإذا بنينا ال!!{structure}
المرضية لـ~$\Gamma$، كفلت هذه الزيادة لكل
\iftag{prvEx}
{جملة وجودية صادقة شاهدًا}
{جملة كلية كاذبة مثالًا ناقضًا}
من بين الرموز الثابتة الجديدة.
\end{explain}

\begin{prop}
\ollabel{prop:lang-exp}
إذا كانت $\Gamma$ متسقة في $\Lang L$، وكانت $\Lang L'$ ناتجة من
$\Lang L$ بإضافة مجموعة !!a{denumerable} من ال!!{constant}s الجديدة
$\Obj d_0$ و$\Obj d_1$ و\dots، فإن $\Gamma$ متسقة في~$\Lang L'$.
\end{prop}

\begin{defn}[المجموعة المشبعة]
المجموعة $\Gamma$ من !!{formula}s اللغة $\Lang {L}$ تسمى
\emph{مشبعة} إذا وفقط إذا كان لكل !!{formula}~$!A(x) \in \Frm[L]$
ليس لها إلا !!{variable} حر واحد~$x$، !!a{constant}~$c \in \Lang{L}$ يحقق
\iftag{prvEx}
      {$\lexists[x][!A(x)] \lif !A(c) \in \Gamma$}
      {$\lnot\lforall[x][!A(x)] \lif \lnot !A(c) \in \Gamma$}.
\end{defn}

ولأجل البرهان الآتي نضع التعريف التالي.

\begin{defn}
\ollabel{defn:henkin-exp}
لتكن $\Lang L'$ اللغة الموصوفة في \olref{prop:lang-exp}. نختار تعدادًا
$!A_0(x_0)$ و$!A_1(x_1)$ و\dots يستوعب ال!!{formula}s~$!A_i(x_i)$
في~$\Lang L'$ التي لها متغير حر واحد ($x_i$). ثم نبني
ال!!{sentence}s~$!D_n$ استقراءً على~$n$.

فنختار $c_0$ أول !!{constant} في قائمة $\Obj d_i$ المضافة
إلى~$\Lang{L}$ مما لا يقع في~$!A_0(x_0)$. وإذا سبق تعريف $!D_0$
و\dots و~$!D_{n-1}$، اخترنا $c_n$ أول ال!!{constant}s الجديدة
~$\Obj d_i$ مما لا يقع في $!D_0$ و\dots و~$!D_{n-1}$، ولا
في~$!A_n(x_n)$.

ونعرّف $!D_{n}$ بأنها ال!!{sentence} \iftag{prvEx}
{$\lexists[x_{n}][!A_{n}(x_{n})] \lif
  !A_{n}(c_{n})$}{$\lnot\lforall[x_{n}][!A_{n}(x_{n})] \lif \lnot
  !A_{n}(c_{n})$}.
\end{defn}

\begin{lem}
\ollabel{lem:henkin}
يمكن توسيع كل مجموعة متسقة~$\Gamma$ إلى مجموعة متسقة مشبعة~$\Gamma'$.
\end{lem}

\begin{proof}
ابدأ بمجموعة جمل متسقة~$\Gamma$ في لغة~$\Lang{L}$، وزد اللغة مجموعة
!!a{denumerable} من ال!!{constant}s الجديدة، ولتكن اللغة الناتجة
~$\Lang{L'}$. تضمن \olref{prop:lang-exp} بقاء اتساق $\Gamma$ بعد
الزيادة. وخذ $!D_i$ من \olref{defn:henkin-exp}، وعرّف
\begin{align*}
\Gamma_0 & = \Gamma \\
\Gamma_{n+1} & = \Gamma_n \cup \{!D_n \}
\end{align*}
فيكون $\Gamma_{n+1} = \Gamma \cup \{ !D_0, \dots, !D_n \}$. ثم ضع
$\Gamma' = \bigcup_{n} \Gamma_n$؛ وإشباع $\Gamma'$ بيّن من بنائها.

إن عدم اتساق $\Gamma'$ يستلزم عدم اتساق $\Gamma_n$ عند بعض $n$
(بيّن وجه ذلك تمرينًا). فإثبات اتساق $\Gamma'$ مردود إلى الاستقراء
على~$n$ لإثبات اتساق كل~$\Gamma_n$.

أما الأساس، فاتساق $\Gamma_0 = \Gamma$ هو الفرض نفسه. وأما الانتقال،
فلنفرض اتساق $\Gamma_{n}$ وعدم اتساق
$\Gamma_{n+1} = \Gamma_n \cup \{!D_n\}$. وقد عرّفنا $!D_n$~بأنها
\iftag{prvEx}
{$\lexists[x_{n}][!A_{n}(x_n)] \lif !A_{n}(c_{n})$}
{$\lnot\lforall[x_{n}][!A_{n}(x_n)] \lif \lnot !A_{n}(c_{n})$},
وصيغة $!A_n(x_n)$ هي !!a{formula} من $\Lang{L'}$ متغيرها الحر
الوحيد~$x_n$. وطريقة اختيار~$c_n$ في
\olref{defn:henkin-exp} تكفل أن $c_n$ لا يقع في~$!A_n(x_n)$ ولا في~$\Gamma_n$.

فمن عدم اتساق $\Gamma_n \cup \{!D_n\}$ يلزم $\Gamma_n
\Proves \lnot !D_n$، ويلزم من ذلك هذان الحكمان:
\iftag{prvEx}{
\[
\Gamma_n \Proves \lexists[x_n][!A_n(x_n)]
\qquad
\Gamma_n \Proves \lnot !A_n(c_n)
\]}{
\[
\Gamma_n \Proves \lnot\lforall[x_n][!A_n(x_n)]
\qquad
\Gamma_n \Proves !A_n(c_n)
\]}
ولأن $c_n$ لا يقع في $\Gamma_n$ ولا في~$!A_n(x_n)$، يصح تطبيق
\tagrefs{prfAX/{fol:axd:qpr:thm:strong-generalization},prfSC/{fol:seq:qpr:thm:strong-generalization},prfND/{fol:ntd:qpr:thm:strong-generalization},prfTab/{fol:tab:qpr:thm:strong-generalization}}.
ومن \iftag{prvEx}
{$\Gamma_n \Proves \lnot !A_n(c_n)$}
{$\Gamma_n \Proves !A_n(c_n)$},
نحصل على
\iftag{prvEx}
{$\Gamma_n \Proves \lforall[x_n][\lnot !A_n(x_n)]$}
{$\Gamma_n \Proves \lforall[x_n][!A_n(x_n)]$}.
فيجتمع الحكمان
\iftag{prvEx}
{$\Gamma_n \Proves \lexists[x_n][!A_n(x_n)]$ و
$\Gamma_n \Proves \lforall[x_n][\lnot !A_n(x_n)]$}
{$\Gamma_n \Proves \lnot\lforall[x_n][!A_n(x_n)]$ و
$\Gamma_n \Proves \lforall[x_n][!A_n(x_n)]$},
ومن اجتماعهما يلزم عدم اتساق $\Gamma_n$ نفسها.
\iftag{prvEx}
{(لاحظ أن
\iftag{prvAll}
{$\lforall[x_n][\lnot !A_n(x_n)] \Proves
  \lnot\lexists[x_n][!A_n(x_n)]$}
{$\lforall[x_n][\lnot !A_n(x_n)]$ تُعرّف بأنها
  $\lnot\lexists[x_n][\lnot\lnot !A_n(x_n)]$ و
  $\lnot\lexists[x_n][\lnot\lnot !A_n(x_n)] \Proves
  \lnot\lexists[x_n][!A_n(x_n)]$}.)}{}
وهذا خلاف اتساق $\Gamma_n$ المفروض. فثبت اتساق
$\Gamma_n \cup \{ !D_n\}$، وتم الاستقراء.
\end{proof}

\begin{explain}
ننتقل إلى المجموعات التي يجتمع فيها الاتساق و\emph{الاكتمال} والإشباع.
سنثبت أن المجموعة من هذا القبيل
\iftag{prvAll}{تضم !!{sentence} مكممة كليًا إذا وفقط إذا ضمت
  جميع حالاتها}{}\iftag{defAll,defEx}{}{، وأنها }\iftag{prvEx}{تضم
  !!{sentence} مكممة وجوديًا إذا وفقط إذا ضمت حالة واحدة منها
  على الأقل}{}. بهذه الخاصية نثبت أن ال!!{structure} المنشأة من
مجموعة كاملة متسقة مشبعة ترضي جميع جمل المجموعة المكممة.
\end{explain}

\begin{prop}\ollabel{prop:saturated-instances}
افترض أن $\Gamma$ كاملة ومتسقة ومشبعة.
\begin{tagenumerate}{prvEx,prvAll}
\tagitem{prvEx}{$\lexists[x][!A(x)] \in \Gamma$ إذا وفقط إذا كان
  $!A(t) \in \Gamma$ لحد مغلق واحد~$t$ على الأقل.}{}
\tagitem{prvAll}{$\lforall[x][!A(x)] \in \Gamma$ إذا وفقط إذا كان
  $!A(t) \in \Gamma$ لكل حد مغلق~$t$.}{}
\end{tagenumerate}
\end{prop}

\begin{proof}
  \begin{tagenumerate}{prvEx,prvAll}
  \tagitem{prvEx}{%
    \iftag{probEx}{تمرين.}{أما إذا كان $\lexists[x][!A(x)]
      \in \Gamma$، فإن إشباع $\Gamma$ يضمن
      $(\lexists[x][!A(x)] \lif !A(c)) \in \Gamma$ لبعض
      ال!!{constant}~$c$. وبحسب
      \tagrefs{prfAX/{fol:axd:ppr:prop:provability-lif},prfSC/{fol:seq:ppr:prop:provability-lif},prfND/{fol:ntd:ppr:prop:provability-lif},prfTab/{fol:tab:ppr:prop:provability-lif}}، البند~(1)،
      و\olref[ccs]{prop:ccs}\olref[ccs]{prop:ccs-prov-in}، نحصل على
      $!A(c) \in \Gamma$.

    وأما العكس فلا يتوقف على الإشباع. فمتى كان $!A(t) \in \Gamma$،
    ثبت $\Gamma \Proves \lexists[x][!A(x)]$ بمقتضى
      \tagrefs{prfAX/{fol:axd:qpr:prop:provability-quantifiers},prfSC/{fol:seq:qpr:prop:provability-quantifiers},prfND/{fol:ntd:qpr:prop:provability-quantifiers},prfTab/{fol:tab:qpr:prop:provability-quantifiers}}، البند~(1). وبحسب
    \olref[ccs]{prop:ccs}\olref[ccs]{prop:ccs-prov-in} نحصل على
    $\lexists[x][!A(x)] \in \Gamma$.}}{}
  \tagitem{prvAll}{%
    \iftag{probAll}{تمرين.}{ليكن $!A(t) \in \Gamma$ لكل حد
      مغلق~$t$. فلو كان
      $\lforall[x][!A(x)] \notin \Gamma$، لأوجب اكتمال $\Gamma$
      أن $\lnot\lforall[x][!A(x)] \in \Gamma$. وبالإشباع نحصل على
      \iftag{prvEx}{$(\lexists[x][\lnot !A(x)] \lif \lnot !A(c)) \in
        \Gamma$}{$(\lnot\lforall[x][!A(x)] \lif \lnot !A(c)) \in
        \Gamma$} لبعض ال!!{constant}~$c$. ثم إن $c$
      حد مغلق، فيقتضي الفرض $!A(c) \in \Gamma$؛ وحينئذ لا تكون $\Gamma$
      متسقة. (تمرين: أقم ال!!{derivation} المبين أن
      \[
      \lnot \lforall[x][!A(x)],
      \iftag{prvEx}{\lexists[x][\lnot !A(x)] \lif \lnot
        !A(c)}{\lnot\lforall[x][!A(x)] \lif \lnot
        !A(c)}, !A(c)
      \]
      غير متسقة.)

      وأما الاتجاه المقابل فلا يفتقر إلى الإشباع. فإذا كان
      $\lforall[x][!A(x)] \in \Gamma$، لزم $\Gamma \Proves !A(t)$
      بمقتضى
      \tagrefs{prfAX/{fol:axd:qpr:prop:provability-quantifiers},prfSC/{fol:seq:qpr:prop:provability-quantifiers},prfND/{fol:ntd:qpr:prop:provability-quantifiers},prfTab/{fol:tab:qpr:prop:provability-quantifiers}},
      البند~(2). ونحصل على $!A(t) \in \Gamma$ بحسب
      \olref[ccs]{prop:ccs}.}}{}
  \end{tagenumerate}
\end{proof}

\end{document}
